\documentclass[11pt,a4paper]{article}
\usepackage[margin=2.5cm]{geometry}
\usepackage{amsmath,amssymb}
\usepackage{graphicx}
\usepackage{tcolorbox}
\usepackage{listings}
\usepackage{xcolor}
\usepackage{hyperref}
\usepackage{tikz}
\usetikzlibrary{shapes,arrows,positioning,fit,backgrounds}

% Color definitions
\definecolor{mlblue}{RGB}{0,102,204}
\definecolor{mlpurple}{RGB}{51,51,178}
\definecolor{mlorange}{RGB}{255,127,14}
\definecolor{mlgreen}{RGB}{44,160,44}

% Key concept box
\newtcolorbox{keybox}[1]{
  colback=mlblue!5!white,
  colframe=mlblue,
  fonttitle=\bfseries,
  title=#1
}

% Example box
\newtcolorbox{examplebox}[1]{
  colback=mlgreen!5!white,
  colframe=mlgreen,
  fonttitle=\bfseries,
  title=#1
}

% Warning box
\newtcolorbox{warningbox}[1]{
  colback=mlorange!5!white,
  colframe=mlorange,
  fonttitle=\bfseries,
  title=#1
}

% Best practice box
\newtcolorbox{bestpracticebox}[1]{
  colback=mlpurple!5!white,
  colframe=mlpurple,
  fonttitle=\bfseries,
  title=#1
}

% Code listing style
\lstset{
  basicstyle=\ttfamily\small,
  backgroundcolor=\color{gray!10},
  frame=single,
  numbers=left,
  numberstyle=\tiny\color{gray},
  breaklines=true,
  captionpos=b
}

% Solidity style
\lstdefinelanguage{Solidity}{
  keywords={contract, function, public, private, view, pure, payable, returns, address, uint, uint256, bool, string, mapping, struct, event, modifier, require, assert, revert, if, else, for, while, return, emit, constructor, memory, storage, calldata, interface, override, virtual},
  comment=[l]{//},
  morecomment=[s]{/*}{*/},
  morestring=[b]",
  sensitive=true
}

\title{\textbf{Lecture Notes: Lesson 4}\\ERC-20 Token Creation}
\author{Cryptocurrency Course}
\date{}

\begin{document}

\maketitle

\tableofcontents
\newpage

\section{Introduction to Token Standards}

\subsection{What Are Token Standards?}

Token standards define common interfaces for smart contracts, enabling interoperability across the Ethereum ecosystem.

\begin{keybox}{Benefits of Standards}
Standardization enables:
\begin{itemize}
  \item \textbf{Wallet Integration}: Any wallet supporting the standard works with any token
  \item \textbf{Exchange Listing}: Exchanges implement standard once for all tokens
  \item \textbf{DApp Compatibility}: DApps interact with tokens uniformly
  \item \textbf{Composability}: Protocols can build on each other
  \item \textbf{Reduced Development}: Leverage existing tools and libraries
\end{itemize}
\end{keybox}

\subsection{Common Token Standards}

\begin{center}
\begin{tabular}{|l|p{10cm}|}
\hline
\textbf{Standard} & \textbf{Purpose} \\
\hline
ERC-20 & Fungible tokens (currencies, utility tokens) \\
ERC-721 & Non-fungible tokens (NFTs, unique items) \\
ERC-1155 & Multi-token (fungible + non-fungible in one contract) \\
ERC-777 & Advanced fungible tokens with hooks \\
ERC-4626 & Tokenized vaults (yield-bearing tokens) \\
\hline
\end{tabular}
\end{center}

\section{ERC-20 Token Standard}

\subsection{Historical Context}

ERC-20 was proposed by Fabian Vogelsteller in November 2015 as Ethereum Request for Comments \#20. It became the de facto standard for fungible tokens.

\textbf{Adoption}: Thousands of tokens use ERC-20, including:
\begin{itemize}
  \item USDT (Tether) - Stablecoin
  \item LINK (Chainlink) - Oracle network token
  \item UNI (Uniswap) - DEX governance token
  \item DAI (MakerDAO) - Decentralized stablecoin
\end{itemize}

\subsection{Interface Specification}

The ERC-20 standard defines 6 mandatory functions and 2 events.

\subsubsection{Required Functions}

\begin{examplebox}{ERC-20 Interface}
\begin{lstlisting}[language=Solidity]
// SPDX-License-Identifier: MIT
pragma solidity ^0.8.0;

interface IERC20 {
    // Returns total token supply
    function totalSupply() external view returns (uint256);

    // Returns balance of an account
    function balanceOf(address account) external view returns (uint256);

    // Transfers tokens from caller to recipient
    function transfer(address recipient, uint256 amount)
        external returns (bool);

    // Returns remaining tokens that spender can transfer
    // on behalf of owner
    function allowance(address owner, address spender)
        external view returns (uint256);

    // Sets amount as the allowance of spender over caller's tokens
    function approve(address spender, uint256 amount)
        external returns (bool);

    // Transfers tokens from sender to recipient using allowance
    function transferFrom(address sender, address recipient, uint256 amount)
        external returns (bool);

    // Emitted when tokens are transferred
    event Transfer(address indexed from, address indexed to, uint256 value);

    // Emitted when allowance is set
    event Approval(address indexed owner, address indexed spender, uint256 value);
}
\end{lstlisting}
\end{examplebox}

\subsubsection{Optional Metadata Functions}

Not required but widely implemented:

\begin{lstlisting}[language=Solidity]
function name() public view returns (string memory);
function symbol() public view returns (string memory);
function decimals() public view returns (uint8);
\end{lstlisting}

\begin{keybox}{Decimals Explanation}
\textbf{decimals} specifies display precision, not blockchain storage.

Example: Token with 18 decimals
\begin{itemize}
  \item Blockchain stores: \texttt{1000000000000000000} (1 with 18 zeros)
  \item User sees: \texttt{1.0} tokens
  \item Formula: $\text{Display} = \text{Storage} / 10^{\text{decimals}}$
\end{itemize}

Common values: 18 (matches ETH), 6 (USDC), 8 (Bitcoin-style)
\end{keybox}

\subsection{Core Mechanics}

\subsubsection{Token Transfer Flow}

\textbf{Direct Transfer} (\texttt{transfer}):
\begin{enumerate}
  \item User calls \texttt{transfer(recipient, amount)}
  \item Contract checks: \texttt{balanceOf[msg.sender] >= amount}
  \item If true: \texttt{balanceOf[msg.sender] -= amount}
  \item And: \texttt{balanceOf[recipient] += amount}
  \item Emit \texttt{Transfer(msg.sender, recipient, amount)}
\end{enumerate}

\textbf{Delegated Transfer} (\texttt{approve} + \texttt{transferFrom}):
\begin{enumerate}
  \item Owner calls \texttt{approve(spender, amount)}
  \item Sets: \texttt{allowance[owner][spender] = amount}
  \item Emit \texttt{Approval(owner, spender, amount)}
  \item Later, spender calls \texttt{transferFrom(owner, recipient, amount)}
  \item Contract checks: \texttt{allowance[owner][spender] >= amount}
  \item If true: Reduce allowance, transfer tokens
\end{enumerate}

\begin{examplebox}{Approve-TransferFrom Use Case}
\textbf{Scenario}: Using Uniswap DEX to swap tokens

\begin{enumerate}
  \item User approves Uniswap contract: \texttt{token.approve(uniswap, 1000)}
  \item User initiates swap on Uniswap UI
  \item Uniswap contract calls: \texttt{token.transferFrom(user, uniswap, 1000)}
  \item Swap executes without requiring user's private key
\end{enumerate}

This pattern enables smart contracts to interact with user tokens safely.
\end{examplebox}

\subsubsection{State Variables}

Minimal ERC-20 implementation requires:

\begin{lstlisting}[language=Solidity]
mapping(address => uint256) private _balances;
mapping(address => mapping(address => uint256)) private _allowances;
uint256 private _totalSupply;
\end{lstlisting}

\textbf{Storage Structure}:
\begin{itemize}
  \item \texttt{\_balances}: Maps each address to its token balance
  \item \texttt{\_allowances}: Nested mapping for approved spenders
  \item \texttt{\_totalSupply}: Sum of all tokens in existence
\end{itemize}

\section{Implementation with OpenZeppelin}

\subsection{Why Use OpenZeppelin?}

OpenZeppelin Contracts is the gold standard library for secure smart contract development.

\textbf{Advantages}:
\begin{itemize}
  \item Audited by multiple security firms
  \item Battle-tested (used by thousands of projects)
  \item Modular and extensible
  \item Well-documented
  \item Regularly updated
\end{itemize}

\subsection{Basic ERC-20 Token}

\begin{examplebox}{Simple Token Implementation}
\begin{lstlisting}[language=Solidity]
// SPDX-License-Identifier: MIT
pragma solidity ^0.8.0;

import "@openzeppelin/contracts/token/ERC20/ERC20.sol";

contract MyToken is ERC20 {
    constructor(uint256 initialSupply) ERC20("MyToken", "MTK") {
        _mint(msg.sender, initialSupply * 10**decimals());
    }
}
\end{lstlisting}

\textbf{Deployment}:
\begin{lstlisting}[language=JavaScript]
// scripts/deploy.js
const initialSupply = 1000000; // 1 million tokens
const MyToken = await ethers.getContractFactory("MyToken");
const token = await MyToken.deploy(initialSupply);
await token.deployed();
console.log("Token deployed to:", token.address);
\end{lstlisting}

This 7-line contract provides:
\begin{itemize}
  \item Full ERC-20 compliance
  \item Minting of initial supply
  \item All transfer and allowance functions
  \item Safe arithmetic (no overflow/underflow)
\end{itemize}
\end{examplebox}

\subsection{Advanced Features}

\subsubsection{Mintable Token}

Allow authorized addresses to create new tokens:

\begin{lstlisting}[language=Solidity]
import "@openzeppelin/contracts/token/ERC20/ERC20.sol";
import "@openzeppelin/contracts/access/Ownable.sol";

contract MintableToken is ERC20, Ownable {
    constructor() ERC20("MintableToken", "MINT") {}

    function mint(address to, uint256 amount) public onlyOwner {
        _mint(to, amount);
    }
}
\end{lstlisting}

\subsubsection{Burnable Token}

Allow token holders to destroy their tokens:

\begin{lstlisting}[language=Solidity]
import "@openzeppelin/contracts/token/ERC20/ERC20.sol";
import "@openzeppelin/contracts/token/ERC20/extensions/ERC20Burnable.sol";

contract BurnableToken is ERC20, ERC20Burnable {
    constructor(uint256 initialSupply) ERC20("BurnableToken", "BURN") {
        _mint(msg.sender, initialSupply);
    }
}
\end{lstlisting}

Usage:
\begin{lstlisting}[language=Solidity]
// Burn your own tokens
token.burn(1000);

// Burn tokens you have allowance for
token.burnFrom(otherAddress, 500);
\end{lstlisting}

\subsubsection{Capped Supply}

Enforce maximum total supply:

\begin{lstlisting}[language=Solidity]
import "@openzeppelin/contracts/token/ERC20/ERC20.sol";
import "@openzeppelin/contracts/token/ERC20/extensions/ERC20Capped.sol";
import "@openzeppelin/contracts/access/Ownable.sol";

contract CappedToken is ERC20, ERC20Capped, Ownable {
    constructor(uint256 cap) ERC20("CappedToken", "CAP")
        ERC20Capped(cap * 10**decimals()) {}

    function mint(address to, uint256 amount) public onlyOwner {
        _mint(to, amount);
    }

    function _mint(address account, uint256 amount)
        internal virtual override(ERC20, ERC20Capped) {
        super._mint(account, amount);
    }
}
\end{lstlisting}

\subsection{Comprehensive Token Example}

\begin{examplebox}{Full-Featured Token}
\begin{lstlisting}[language=Solidity]
// SPDX-License-Identifier: MIT
pragma solidity ^0.8.0;

import "@openzeppelin/contracts/token/ERC20/ERC20.sol";
import "@openzeppelin/contracts/token/ERC20/extensions/ERC20Burnable.sol";
import "@openzeppelin/contracts/token/ERC20/extensions/ERC20Capped.sol";
import "@openzeppelin/contracts/access/Ownable.sol";
import "@openzeppelin/contracts/security/Pausable.sol";

contract ComprehensiveToken is ERC20, ERC20Burnable,
    ERC20Capped, Ownable, Pausable {

    // Events
    event TokensMinted(address indexed to, uint256 amount);
    event TokensBurned(address indexed from, uint256 amount);

    constructor(
        string memory name,
        string memory symbol,
        uint256 cap,
        uint256 initialSupply
    ) ERC20(name, symbol) ERC20Capped(cap * 10**decimals()) {
        require(initialSupply <= cap, "Initial supply exceeds cap");
        _mint(msg.sender, initialSupply * 10**decimals());
    }

    // Minting function (only owner)
    function mint(address to, uint256 amount) public onlyOwner {
        _mint(to, amount);
        emit TokensMinted(to, amount);
    }

    // Enhanced burn with event
    function burn(uint256 amount) public override {
        super.burn(amount);
        emit TokensBurned(msg.sender, amount);
    }

    // Pause/unpause transfers
    function pause() public onlyOwner {
        _pause();
    }

    function unpause() public onlyOwner {
        _unpause();
    }

    // Override required by Solidity
    function _mint(address account, uint256 amount)
        internal virtual override(ERC20, ERC20Capped) {
        super._mint(account, amount);
    }

    // Ensure transfers respect pause state
    function _beforeTokenTransfer(
        address from,
        address to,
        uint256 amount
    ) internal override whenNotPaused {
        super._beforeTokenTransfer(from, to, amount);
    }
}
\end{lstlisting}

\textbf{Features}:
\begin{itemize}
  \item ERC-20 compliant
  \item Burnable (reduce supply)
  \item Capped (maximum supply)
  \item Mintable (owner only)
  \item Pausable (emergency stop)
  \item Ownable (access control)
\end{itemize}
\end{examplebox}

\section{Token Economics and Supply Models}

\subsection{Supply Mechanisms}

\subsubsection{Fixed Supply}

Total supply created at deployment, never changes.

\textbf{Example}: Bitcoin (21M cap), Wrapped Bitcoin (WBTC)

\textbf{Implementation}:
\begin{lstlisting}[language=Solidity]
constructor() ERC20("FixedToken", "FIX") {
    _mint(msg.sender, 21_000_000 * 10**decimals());
}
\end{lstlisting}

\textbf{Characteristics}:
\begin{itemize}
  \item Deflationary pressure if tokens are lost
  \item Scarcity-driven value
  \item Predictable supply schedule
\end{itemize}

\subsubsection{Inflationary Supply}

New tokens continuously minted.

\textbf{Example}: Ethereum (before EIP-1559 burn), staking rewards

\textbf{Implementation}:
\begin{lstlisting}[language=Solidity]
uint256 public constant ANNUAL_INFLATION_RATE = 2; // 2% per year

function mintInflation() public onlyOwner {
    uint256 timeSinceLastMint = block.timestamp - lastMintTimestamp;
    uint256 annualSeconds = 365 days;
    uint256 inflationAmount = (totalSupply() * ANNUAL_INFLATION_RATE
        * timeSinceLastMint) / (100 * annualSeconds);

    _mint(treasury, inflationAmount);
    lastMintTimestamp = block.timestamp;
}
\end{lstlisting}

\textbf{Characteristics}:
\begin{itemize}
  \item Rewards participation (staking, mining)
  \item Can decrease token value over time
  \item Funds development or ecosystem growth
\end{itemize}

\subsubsection{Elastic Supply}

Supply adjusts to maintain price target (rebase tokens).

\textbf{Example}: Ampleforth (AMPL)

\textbf{Mechanism}: If price > \$1.05, increase all balances proportionally. If price < \$0.95, decrease balances.

\begin{warningbox}{Rebase Tokens}
Elastic supply tokens are NOT standard ERC-20 compatible:
\begin{itemize}
  \item Balance changes without transfers
  \item Breaks DeFi integrations expecting constant balances
  \item Requires special handling by exchanges/wallets
\end{itemize}
\end{warningbox}

\subsection{Distribution Strategies}

\subsubsection{Initial Coin Offering (ICO)}

\begin{itemize}
  \item Sell tokens to public before launch
  \item Raised \$20B+ in 2017-2018 peak
  \item Often unregulated, high fraud risk
  \item Regulatory crackdowns in most jurisdictions
\end{itemize}

\subsubsection{Airdrop}

\begin{itemize}
  \item Free distribution to specific addresses
  \item Reward early users or protocol participants
  \item Marketing and community building
\end{itemize}

\begin{examplebox}{Airdrop Implementation}
\begin{lstlisting}[language=Solidity]
contract Airdrop {
    IERC20 public token;
    mapping(address => bool) public claimed;

    function claimAirdrop() public {
        require(!claimed[msg.sender], "Already claimed");
        require(isEligible(msg.sender), "Not eligible");

        claimed[msg.sender] = true;
        token.transfer(msg.sender, 100 * 10**18); // 100 tokens
    }

    function isEligible(address user) internal view returns (bool) {
        // Example: Must have interacted before block X
        // Implementation depends on criteria
        return true;
    }
}
\end{lstlisting}
\end{examplebox}

\subsubsection{Liquidity Mining}

\begin{itemize}
  \item Reward users who provide liquidity to DEXs
  \item Tokens distributed over time based on participation
  \item Popularized by Compound, Uniswap, etc.
\end{itemize}

\subsubsection{Vesting}

Lock tokens and release gradually over time.

\begin{examplebox}{Simple Vesting}
\begin{lstlisting}[language=Solidity]
contract TokenVesting {
    IERC20 public token;
    address public beneficiary;
    uint256 public cliff;        // When vesting starts
    uint256 public duration;     // Vesting period
    uint256 public released;

    constructor(
        IERC20 _token,
        address _beneficiary,
        uint256 _cliff,
        uint256 _duration
    ) {
        token = _token;
        beneficiary = _beneficiary;
        cliff = block.timestamp + _cliff;
        duration = _duration;
    }

    function release() public {
        require(block.timestamp >= cliff, "Cliff not reached");

        uint256 vested = vestedAmount();
        uint256 releasable = vested - released;
        require(releasable > 0, "No tokens to release");

        released += releasable;
        token.transfer(beneficiary, releasable);
    }

    function vestedAmount() public view returns (uint256) {
        uint256 totalBalance = token.balanceOf(address(this)) + released;

        if (block.timestamp < cliff) {
            return 0;
        } else if (block.timestamp >= cliff + duration) {
            return totalBalance;
        } else {
            return (totalBalance * (block.timestamp - cliff)) / duration;
        }
    }
}
\end{lstlisting}
\end{examplebox}

\subsection{Tokenomics Design Considerations}

\begin{enumerate}
  \item \textbf{Utility}: What does the token do?
    \begin{itemize}
      \item Governance (voting rights)
      \item Access (pay for services)
      \item Staking (earn rewards, secure network)
      \item Dividends (share revenue)
    \end{itemize}

  \item \textbf{Supply Schedule}: When are tokens released?
    \begin{itemize}
      \item Avoid dumping large amounts at once
      \item Align incentives with long-term growth
      \item Consider market absorption capacity
    \end{itemize}

  \item \textbf{Allocation}: Who gets tokens?
    \begin{itemize}
      \item Team (10-20\%, vested)
      \item Investors (10-30\%, vested)
      \item Community/Ecosystem (40-60\%)
      \item Treasury/Reserve (10-20\%)
    \end{itemize}

  \item \textbf{Value Accrual}: How does token capture value?
    \begin{itemize}
      \item Token burns (reduce supply)
      \item Buybacks (market support)
      \item Staking yield (incentivize holding)
      \item Revenue sharing (cashflow to holders)
    \end{itemize}
\end{enumerate}

\section{Security Considerations}

\subsection{Common Vulnerabilities}

\subsubsection{Integer Overflow/Underflow}

\textbf{Pre-Solidity 0.8.0 Issue}:
\begin{lstlisting}[language=Solidity]
// Vulnerable (Solidity < 0.8.0)
uint256 balance = 0;
balance -= 1; // Underflows to 2^256 - 1
\end{lstlisting}

\textbf{Solution}: Solidity 0.8+ has built-in overflow checking. For older versions, use SafeMath library.

\subsubsection{Approval Race Condition}

\textbf{Problem}: Changing approval from N to M can be exploited.

\textbf{Attack Scenario}:
\begin{enumerate}
  \item Alice approves Bob for 100 tokens
  \item Alice decides to change to 50, sends \texttt{approve(Bob, 50)}
  \item Bob sees pending transaction, front-runs with \texttt{transferFrom(Alice, Bob, 100)}
  \item Bob's transaction executes first (uses old 100 allowance)
  \item Alice's transaction executes (sets allowance to 50)
  \item Bob calls \texttt{transferFrom(Alice, Bob, 50)} again
  \item Total stolen: 150 tokens
\end{enumerate}

\textbf{Mitigation}:
\begin{lstlisting}[language=Solidity]
// Set to 0 first, then set to new value
token.approve(spender, 0);
token.approve(spender, newAmount);

// Or use increaseAllowance/decreaseAllowance
token.increaseAllowance(spender, 50);
token.decreaseAllowance(spender, 30);
\end{lstlisting}

\subsubsection{Unlimited Approvals}

\textbf{Common Pattern}: \texttt{approve(spender, type(uint256).max)}

\textbf{Risk}: If spender contract is compromised, all tokens can be stolen.

\bestpracticebox{Approval Best Practices}
\begin{itemize}
  \item Approve only what's needed for immediate transaction
  \item Revoke approvals after use
  \item Use tools like \url{https://revoke.cash} to audit/revoke approvals
  \item Consider using permit (EIP-2612) for gasless approvals
\end{itemize}
\end{bestpracticebox}

\subsection{Audit and Testing}

\subsubsection{Testing Framework}

\begin{examplebox}{Hardhat Test Example}
\begin{lstlisting}[language=JavaScript]
const { expect } = require("chai");

describe("MyToken", function() {
  let token, owner, addr1, addr2;

  beforeEach(async function() {
    [owner, addr1, addr2] = await ethers.getSigners();
    const Token = await ethers.getContractFactory("MyToken");
    token = await Token.deploy(1000000);
  });

  it("Should assign total supply to owner", async function() {
    const ownerBalance = await token.balanceOf(owner.address);
    expect(await token.totalSupply()).to.equal(ownerBalance);
  });

  it("Should transfer tokens correctly", async function() {
    await token.transfer(addr1.address, 50);
    expect(await token.balanceOf(addr1.address)).to.equal(50);
  });

  it("Should fail if sender doesn't have enough tokens", async function() {
    const initialOwnerBalance = await token.balanceOf(owner.address);
    await expect(
      token.connect(addr1).transfer(owner.address, 1)
    ).to.be.revertedWith("ERC20: transfer amount exceeds balance");
  });

  it("Should handle allowance correctly", async function() {
    await token.approve(addr1.address, 100);
    expect(await token.allowance(owner.address, addr1.address))
      .to.equal(100);

    await token.connect(addr1).transferFrom(
      owner.address, addr2.address, 50
    );
    expect(await token.balanceOf(addr2.address)).to.equal(50);
    expect(await token.allowance(owner.address, addr1.address))
      .to.equal(50);
  });
});
\end{lstlisting}
\end{examplebox}

\subsubsection{Security Audit Checklist}

\begin{enumerate}
  \item \textbf{Access Control}
    \begin{itemize}
      \item Only authorized addresses can mint
      \item Owner cannot steal user tokens
      \item Ownership transfer is secure
    \end{itemize}

  \item \textbf{Arithmetic}
    \begin{itemize}
      \item No overflow/underflow possible
      \item Rounding errors minimized
      \item Division by zero prevented
    \end{itemize}

  \item \textbf{External Calls}
    \begin{itemize}
      \item Reentrancy protection where needed
      \item Check return values
      \item Gas limits considered
    \end{itemize}

  \item \textbf{Token Mechanics}
    \begin{itemize}
      \item Transfer logic correct
      \item Allowance handling secure
      \item Events emitted properly
      \item Decimal handling consistent
    \end{itemize}

  \item \textbf{Deployment}
    \begin{itemize}
      \item Constructor parameters validated
      \item Initial supply distributed correctly
      \item Contract ownership set properly
    \end{itemize}
\end{enumerate}

\section{Deployment and Verification}

\subsection{Deployment Process}

\begin{examplebox}{Complete Deployment Script}
\begin{lstlisting}[language=JavaScript]
const hre = require("hardhat");

async function main() {
  // Deployment parameters
  const NAME = "My Project Token";
  const SYMBOL = "MPT";
  const CAP = 100_000_000;        // 100 million cap
  const INITIAL_SUPPLY = 10_000_000; // 10 million initial

  console.log("Deploying token...");

  const Token = await hre.ethers.getContractFactory(
    "ComprehensiveToken"
  );
  const token = await Token.deploy(
    NAME,
    SYMBOL,
    CAP,
    INITIAL_SUPPLY
  );

  await token.deployed();

  console.log("Token deployed to:", token.address);
  console.log("Name:", await token.name());
  console.log("Symbol:", await token.symbol());
  console.log("Decimals:", await token.decimals());
  console.log("Total Supply:", ethers.utils.formatEther(
    await token.totalSupply()
  ));
  console.log("Cap:", ethers.utils.formatEther(await token.cap()));

  // Wait for 5 confirmations before verifying
  console.log("Waiting for block confirmations...");
  await token.deployTransaction.wait(5);

  // Verify on Etherscan
  console.log("Verifying contract...");
  await hre.run("verify:verify", {
    address: token.address,
    constructorArguments: [NAME, SYMBOL, CAP, INITIAL_SUPPLY],
  });

  console.log("Deployment complete!");
}

main().catch((error) => {
  console.error(error);
  process.exitCode = 1;
});
\end{lstlisting}

Execution:
\begin{verbatim}
npx hardhat run scripts/deploy.js --network sepolia
\end{verbatim}
\end{examplebox}

\subsection{Contract Verification}

Contract verification publishes source code on block explorers, enabling:
\begin{itemize}
  \item Users to read contract code
  \item Direct interaction via explorer UI
  \item Trust and transparency
  \item Automated security checks
\end{itemize}

\textbf{hardhat.config.js}:
\begin{lstlisting}[language=JavaScript]
require("@nomiclabs/hardhat-etherscan");

module.exports = {
  solidity: "0.8.20",
  networks: {
    sepolia: {
      url: process.env.SEPOLIA_RPC_URL,
      accounts: [process.env.PRIVATE_KEY]
    }
  },
  etherscan: {
    apiKey: process.env.ETHERSCAN_API_KEY
  }
};
\end{lstlisting}

\section{Study Questions}

\subsection{Conceptual Questions}

\begin{enumerate}
  \item Explain the difference between \texttt{transfer} and \texttt{transferFrom}.
  \item Why is the approve-transferFrom pattern necessary for DEX interactions?
  \item What are the risks of unlimited token approvals?
  \item How does \texttt{decimals} affect token display vs. storage?
  \item Compare fixed supply vs. inflationary token economics.
\end{enumerate}

\subsection{Implementation Exercises}

\begin{enumerate}
  \item Create an ERC-20 token with 2\% transfer tax that goes to a treasury.
  \item Implement a token with blacklist functionality to freeze malicious accounts.
  \item Build a token that pays dividends to holders from contract balance.
  \item Design a governance token where voting power increases with holding duration.
\end{enumerate}

\subsection{Security Analysis}

\begin{enumerate}
  \item Identify the approval race condition vulnerability and propose mitigation.
  \item Write tests to verify all critical token functions.
  \item Analyze a given token contract for potential security issues.
  \item Design a safe token vesting contract with multiple beneficiaries.
\end{enumerate}

\subsection{Advanced Topics}

\begin{enumerate}
  \item Implement EIP-2612 (permit) for gasless approvals.
  \item Create a token with snapshot functionality for airdrops.
  \item Build a wrapper token (like WETH) that converts ETH to ERC-20.
  \item Design a rebasing token that adjusts supply based on oracle price.
\end{enumerate}

\section{Further Reading}

\begin{itemize}
  \item EIP-20: Token Standard Specification
  \item OpenZeppelin ERC-20 Documentation and Code
  \item EIP-2612: Permit Extension (Gasless Approvals)
  \item EIP-4626: Tokenized Vault Standard
  \item ConsenSys: Token Implementation Best Practices
  \item Trail of Bits: Token Security Checklist
  \item DeFi Developer Roadmap: Token Standards
  \item Ethereum.org: Token Standards Overview
\end{itemize}

\end{document}
